\section{Rangkuman}
\label{sec:rangkuman-bab9}

Pada bab ini telah dibahas:

\begin{itemize}
    \item Kode sumber assembly terdiri atas dua jenis baris:
          instruksi mesin yang menghasilkan opcode dan dieksekusi
          saat program berjalan, serta direktif assembler yang
          hanya berlaku pada saat perakitan dan tidak menghasilkan
          kode mesin \cite{hyde2010}.

    \item \dir{DB} menyediakan satu byte, \dir{DW} dua byte,
          \dir{DD} empat byte, dan \dir{DT} sepuluh byte. Operator
          \dir{DUP} mengulang satu pola, sedangkan penanda \code{?}
          hanya menyediakan ruang tanpa nilai awal.

    \item Bilangan multi-byte disimpan dengan urutan
          \eng{little-endian}, sehingga \lstasm|DW 1234h|
          tersimpan sebagai byte \hex{34} lalu \hex{12}.

    \item \dir{EQU} membuat konstanta tetap yang tidak dapat
          didefinisikan ulang, sedangkan \code{=} membuat nilai
          yang boleh diganti berkali-kali.

    \item Direktif lokasi dan struktur meliputi \dir{ORG},
          \dir{SEGMENT}/\dir{ENDS}, \dir{ASSUME}, dan \dir{END}.
          Direktif \dir{END} yang diberi label sekaligus
          menetapkan titik masuk program.

    \item \dir{.MODEL} memilih model memori: \file{TINY} satu
          segmen, \file{SMALL} satu segmen kode dan satu segmen
          data, sedangkan model yang lebih besar memakai banyak
          segmen.

    \item Assembler bekerja dalam dua lintasan: membangun tabel
          simbol dan menghitung \eng{offset}, lalu menghasilkan
          kode mesin. Berkas \file{.lst} memperlihatkan hasilnya,
          termasuk penanda \code{R} \cite{bodur}. Linker kemudian
          menyelesaikan simbol eksternal, melakukan
          \eng{relocation}, dan merekam alamat akhir pada berkas
          \file{.map} \cite{nasm2024}.

    \item Alur lengkap dari DOS adalah
          \file{tasm /l program.asm} $\rightarrow$
          \file{tlink /m program.obj} $\rightarrow$ menjalankan
          \file{program.exe}.
\end{itemize}

\begin{checklistbox}
Tandai setiap butir yang sudah Anda kuasai.

\begin{itemize}
    \item[\checkbox] Saya dapat menjelaskan mengapa direktif tidak
          menghasilkan kode mesin.
    \item[\checkbox] Saya dapat memilih direktif data yang tepat
          dan menghitung besar memori yang disediakannya.
    \item[\checkbox] Saya dapat menentukan urutan byte
          \eng{little-endian} pada \dir{DW} dan \dir{DD}.
    \item[\checkbox] Saya dapat menjelaskan perbedaan \dir{EQU} dan
          \code{=} beserta akibat definisi ulang.
    \item[\checkbox] Saya dapat menuliskan struktur segmen program
          \file{.COM} dan \file{.EXE}.
    \item[\checkbox] Saya dapat memilih model memori \dir{.MODEL}
          dan menjelaskan pengaruhnya pada jumlah segmen.
    \item[\checkbox] Saya dapat membaca tabel simbol pada berkas
          \file{.lst} dan tabel segmen pada berkas \file{.map}.
    \item[\checkbox] Saya dapat merakit dan mengaitkan program
          dengan TASM dan TLINK tanpa kesalahan.
\end{itemize}

\vspace{0.4em}
Bila masih ada butir yang belum bertanda, ulangi bagian yang
bersesuaian sebelum melanjutkan ke bab berikutnya.
\end{checklistbox}

\begin{infobox}
\textbf{Rujukan bab.} Pembahasan direktif data dan peran assembler
diperkuat oleh \cite{hyde2010} dan \cite{bodur}. Keterangan praktis
mengenai direktif \dir{EQU}, \dir{ORG}, dan \dir{DUP} pada assembler
modern dapat ditelusuri pada \cite{nasm2024}, sedangkan contoh
pemakaiannya pada lingkungan belajar tersedia pada
\cite{emu8086help}.
\end{infobox}
